\section*{अध्याय \ref{ch:b3:cell-cycle-apoptosis} --- कोशिका-चक्र का नियंत्रण और क्रमादेशित कोशिका-मृत्यु}

\begin{solution}{exo:b3:cell-cycle-apoptosis:1}
G1, S, G2, M। G1 की प्रगति और \omterm{def:b3:cell-cycle-apoptosis:checkpoints}{निर्बंधन बिंदु}: \omterm{def:b3:cell-cycle-apoptosis:cdk}{चक्रिन} D--Cdk4/6;
G1/S: \omterm{def:b3:cell-cycle-apoptosis:cdk}{चक्रिन} E--Cdk2; S: \omterm{def:b3:cell-cycle-apoptosis:cdk}{चक्रिन} A--Cdk2; G2/M: \omterm{def:b3:cell-cycle-apoptosis:cdk}{चक्रिन} B--Cdk1।
\omterm{def:b3:cell-cycle-apoptosis:cdk}{पश्चावस्था-प्रवर्तक संकुल} (APC/C) \omterm{def:b3:cell-cycle-apoptosis:cdk}{चक्रिन} B और सेक्युरिन नष्ट करता है और
समसूत्री विभाजन समाप्त करता है।
\end{solution}

\begin{solution}{exo:b3:cell-cycle-apoptosis:2}
हार्टवेल: मुकुलक खमीर के \emph{cdc} उत्परिवर्ती, जो निश्चित अवस्थाओं पर
रुके थे, और Start की संकल्पना। नर्स: विखंडन खमीर का \emph{cdc2}
उस काइनेज़ के रूप में जो समसूत्री विभाजन का समय तय करती है, और उसका
मानव समजात (Cdk1) जो खमीर को उबार लेता है --- संरक्षण। हंट: समुद्री
साही के अंडों में \omterm{def:b3:cell-cycle-apoptosis:cdk}{चक्रिन}, ऐसा प्रोटीन जो चक्र भर संश्लेषित होता है और
हर विभाजन पर नष्ट हो जाता है।
\end{solution}

\begin{solution}{exo:b3:cell-cycle-apoptosis:3}
\omterm{def:b3:cell-cycle-apoptosis:checkpoints}{निर्बंधन बिंदु} (पछेती G1): वृद्धि कारक, पोषक, आकार; \omterm{def:b3:cell-cycle-apoptosis:cdk}{चक्रिन} D--Cdk4/6
और Rb पर काम करता है। G2/M: DNA अक्षत और प्रतिकृत; \omterm{def:b3:dna-repair:ddr}{ATM}/ATR
Cdc25 को रोकते हैं, जिससे Cdk1 फ़ॉस्फ़ोरिलित और निष्क्रिय रहता है। तर्कु
संयोजन: हर गुणसूत्रबिंदु जुड़ा हो; न जुड़े गुणसूत्रबिंदु
Mad2--Cdc20 संकुल बनाते हैं जो APC/C को रोकते हैं।
\end{solution}

\begin{solution}{exo:b3:cell-cycle-apoptosis:4}
\omterm{def:b3:cell-cycle-apoptosis:apoptosis}{एपॉप्टोसिस}: कोशिका सिकुड़ती है, \omterm{def:b3:chromatin-epigenetics:nucleosome}{क्रोमैटिन} संघनित होता है, DNA सीढ़ी में
कटता है, झिल्ली बुलबुलाती है पर बंद रहती है, \omterm{def:b3:cell-cycle-apoptosis:apoptosis}{फ़ॉस्फ़ेटिडिलसेरीन} उजागर
होता है, मुर्दा निगल लिया जाता है, कोई शोथ नहीं। \omterm{def:b3:cell-cycle-apoptosis:apoptosis}{परिगलन}: कोशिका फूलती
है, झिल्ली फटती है, अंतर्वस्तु रिसती है, DNA का यादृच्छिक अपघटन, शोथ।
निष्पादक: \omterm{def:b3:cell-cycle-apoptosis:apoptosis}{कैस्पेज़}, यानी ऐस्पार्टेट के बाद काटती सिस्टीन प्रोटिएज़।
\end{solution}

\begin{solution}{exo:b3:cell-cycle-apoptosis:5}
$5$ से $30$ तक प्रति मिनट $1$ की दर से उठान: \qty{25}{min}।
\qty{2}{min} अर्ध-आयु काल के साथ $30$ से $5$ तक अपघटन: $\log_{2}(30/5)\times 2 =
\qty{5.2}{min}$। आवर्तकाल लगभग \qty{30}{min}, जिसका \qty{17}{\%} समसूत्री विभाजन।
\end{solution}

\begin{solution}{exo:b3:cell-cycle-apoptosis:6}
(क) $C$ कभी $C_{1}$ से नीचे नहीं गिर सकता: उच्च शाखा पर समसूत्री
विभाजन में बंद। (ख) कोई निरोधी फ़ॉस्फ़ोरिलन नहीं: देहली $C_{2}$ नीची हो
जाती है, समसूत्री विभाजन छोटे आकार पर जल्दी शुरू हो जाता है (``वी''
लक्षण-प्ररूप)। (ग) Cdk1 फ़ॉस्फ़ोरिलित बना रहता है: उच्च शाखा तक पहुँचा
ही नहीं जा सकता, और G2 में लंबी कोशिकाओं के साथ विराम। (घ) जैसा (ख)
में, और G2/M \omterm{def:b3:cell-cycle-apoptosis:checkpoints}{जाँच-बिंदु}, जो Wee1/Cdc25 के ज़रिए काम करता है, कोशिका को
अब थाम नहीं सकता।
\end{solution}

\begin{solution}{exo:b3:cell-cycle-apoptosis:7}
S-प्रावस्था कोशिकाद्रव्य में G1 केंद्रक: उसके उद्गम अनुज्ञप्त हैं और
कोशिकाद्रव्य सक्रिय \omterm{def:b3:cell-cycle-apoptosis:cdk}{चक्रिन} E/A--Cdk2 देता है, इसलिए वह तुरंत S में घुस
जाता है। S-प्रावस्था कोशिकाद्रव्य में G2 केंद्रक: उसके उद्गम चल चुके
और फिर अनुज्ञप्त नहीं हुए (अनुज्ञापन-कारक CDK द्वारा नष्ट या निरुद्ध),
इसलिए वह फिर प्रतिकृति नहीं कर सकता; वह तब तक प्रतीक्षा करता है जब तक
साथी G2 तक न पहुँच जाए और दोनों साथ समसूत्री विभाजन में न चले जाएँ।
\end{solution}

\begin{solution}{exo:b3:cell-cycle-apoptosis:8}
न जुड़ा गुणसूत्रबिंदु एक उत्प्रेरक है: वह Mad2 को लगातार उसकी सक्रिय
संरूपता में बदलता रहता है और Cdc20 का ऐसा विसरणशील निरोधक बनाता है जो
कोशिका भर फैल जाता है और APC/C को हर जगह बंद रखता है --- एक ही
गुणसूत्रबिंदु पर्याप्त है। Mad2 के बिना APC/C तभी सक्रिय हो जाता है जब
Cdk1 उसे चालू कर दे, चाहे गुणसूत्र जुड़े हों या नहीं: पश्चावस्था न जुड़े
गुणसूत्रों के साथ शुरू हो जाती है, पुत्रियाँ असुगुणित होती हैं, और जीव
(या कोशिका-वंश) गुणसूत्र-ह्रास से मर जाता है।
\end{solution}

\begin{solution}{exo:b3:cell-cycle-apoptosis:9}
\emph{ced-9} नहीं: जिन कोशिकाओं को जीना चाहिए वे मर जाती हैं ---
बहुत अधिक मृत्यु से भ्रूणीय घातकता। \emph{ced-3} नहीं: उन $131$
मृत्युओं में से कोई नहीं होती, कोशिकाएँ बचकर विभेदित हो जाती हैं, कृमि
जीवक्षम है। दोनों नहीं: कोई मृत्यु नहीं --- यानी \emph{ced-3}
लक्षण-प्ररूप। चूँकि निष्पादक हटाने से रक्षक हटाने का प्रभाव रद्द हो
जाता है, CED-9 ऊपर की ओर काम करता है, यानी CED-4/CED-3 को निष्क्रिय
थामकर; जब वह न हो तो वे सक्रिय रहते हैं, बशर्ते वे स्वयं अनुपस्थित न हों।
\end{solution}

\begin{solution}{exo:b3:cell-cycle-apoptosis:10}
$A = aC$ के साथ $\mathrm{d}C/\mathrm{d}t = k_{s} - k_{d}aC^{2}$, यानी
एक-चर समीकरण जिसकी स्थायी अवस्था $C^{*} = \sqrt{k_{s}/(k_{d}a)}$ है
और वहाँ ढाल $-2k_{d}aC^{*} < 0$: कोई भी विचलन एकदिष्ट रूप से क्षीण हो
जाता है, और कोई अकेला प्रथम-कोटि चर दोलन नहीं कर सकता।
S-आकार वक्र दो स्थिर शाखाएँ देता है जिनके बीच कोई वर्जित क्षेत्र है,
इसलिए अवस्था को कूदना पड़ता है और वह जम नहीं सकती; पुनर्भरण में कोई
स्पष्ट विलंब किसी और रास्ते से दोलन देता --- ब्रेक का बहुत देर से आना,
आगे निकल जाना, और इसी तरह --- जैसा अनेक हॉर्मोनी लयों में होता है।
\end{solution}

\begin{solution}{exo:b3:cell-cycle-apoptosis:11}
हर आता संवेदक किसी रक्षक द्वारा पकड़ लिया जाता है, इसलिए तब तक कुछ नहीं
होता जब तक सारे $N$ रक्षक अधिकृत न हो जाएँ, यानी $t \approx N/r$ पर;
अगले संवेदक मुक्त रहते हैं, Bax/Bak सक्रिय करते हैं, जिनके छिद्र-निर्माण
की दर उनकी संख्या के साथ तेज़ी से चढ़ती है, और सूत्रकणिकाएँ मिनटों में
खुल जाती हैं। कोई प्रबल उद्दीपन समय केवल $r$ के ज़रिए घटाता है; देहली
के ऊपर परिणाम वही रहता है। वेनेटोक्लैक्स रक्षकों की जगह घेर लेती है और
प्रभावी $N$ घटा देती है: समय सिकुड़ जाता है, और कोई तैयार कोशिका ---
पहले से $N$ के पास --- तुरंत मर जाती है।
\end{solution}

\begin{solution}{exo:b3:cell-cycle-apoptosis:12}
अतिरिक्त Bcl-2 \omterm{def:b3:cell-cycle-apoptosis:pathways}{आंतरिक पथ} रोक देता है: जिन B~कोशिकाओं को अपने
वृद्धि-संकेत समाप्त होने पर मर जाना चाहिए वे बच जाती हैं, और समष्टि
मृत्यु की विफलता से बढ़ती है, उसी धीमी दर पर जिस दर पर ऐसी कोशिकाएँ
बनती हैं --- मंद, जब तक कोई दूसरा घाव प्रचुरोद्भवन न जोड़ दे। Rb का
ह्रास \omterm{def:b3:cell-cycle-apoptosis:checkpoints}{निर्बंधन बिंदु} हटा देता है: दृष्टिपटल के पूर्वगामी बिना वृद्धि
कारकों के चक्र में आगे बढ़ते हैं और उतनी ही तेज़ी से विभाजित होते हैं
जितनी इंजन देता है --- उग्र। दोनों अर्बुद दोनों कार्यक्रमों के
अलग-अलग विफल होने के उदाहरण हैं: एक मृत्यु-नियंत्रण और एक
विभाजन-नियंत्रण, जिनमें से किसी एक का भी ह्रास किसी अर्बुद के लिए
पर्याप्त है, और दूसरा तेज़ है।
\end{solution}

\begin{solution}{pb:b3:cell-cycle-apoptosis:1}
\textbf{1.} प्रति दिन $10^{7}\times 3500/5 = 7\times 10^{9}$ कोशिकाएँ,
यानी प्रति सेकंड लगभग \num{80000}।
\textbf{2.} प्रति गर्तिका प्रति दिन $3500/5 = 700$; $22\times 2^{5} = 704$।
\textbf{3.} $80\times 365 \approx \num{29000}$ विभाजन; $\times 0.6
\approx \num{17500}$ उत्परिवर्तन।
\textbf{4.} संक्रमण-कोशिकाएँ दिनों में झाड़ दी जाती हैं और अपने
उत्परिवर्तन साथ ले जाती हैं; स्तंभ कोशिका जीवन भर बनी रहती है, इसलिए
उसका हर उत्परिवर्तन रखा जाता है और हर विभाजन और जोड़ता है --- धीमा
चक्रण उन्हें न्यूनतम करता है।
\textbf{5.} रसांकुर अपने सामान्य संक्रमण के लगभग उन्हीं $5$ दिनों में
खाली हो जाता है; उसे फिर बनाने में उबरने का एक दिन, \qty{12}{h} पर
पाँच विभाजन और रसांकुर पर चढ़ाई लगती है --- चार से पाँच दिन, जिनमें
अवरोध नंगा रहता है: आँत के तेज़ संक्रमण-विभाजन ही उसे विकिरण की अगेती
शिकार बनाते हैं।
\textbf{6.} कोई बंद एपॉप्टोटिक मुर्दा जीवाणुओं से भरे किसी अवकाश में
अपनी अंतर्वस्तु नहीं छोड़ता और कोई शोथ नहीं भड़काता; \omterm{def:b3:cell-cycle-apoptosis:apoptosis}{परिगलन} हर झड़ती
कोशिका पर एंज़ाइम और संकेत बिखेरकर श्लेष्मकला में शोथ कर देता।
\textbf{7.} $(40 - 8)/2 = \qty{16}{h}$।
\textbf{8.} $\log_{2}(40/8)\times \qty{10}{min} = 2.3\times 10 =
\qty{23}{min}$।
\textbf{9.} आवर्तकाल लगभग \qty{16.4}{h}; Cdk1 समय का \qty{2.3}{\%}
सक्रिय।
\textbf{10.} अकेला $k_{s}$ नहीं, बल्कि देहलियों पर रुकने का समय: कायिक
कोशिका \omterm{def:b3:cell-cycle-apoptosis:checkpoints}{निर्बंधन बिंदु} पर वृद्धि कारकों की और G2/M पर अपने DNA की
प्रतीक्षा करती है, और स्तंभ कोशिका संक्रमण-कोशिका से अधिक। मेंढक के
अंडे में कोई \omterm{def:b3:cell-cycle-apoptosis:checkpoints}{जाँच-बिंदु} नहीं और कोशिकाद्रव्य हर चीज़ से भरा है, इसलिए
उसका आवर्तकाल केवल नंगा दोलक तय करता है।
\textbf{11.} सामान्य $C_{2}$ से ऊपर के किसी \omterm{def:b3:cell-cycle-apoptosis:cdk}{चक्रिन} स्तर पर निम्न शाखा
पर: \omterm{def:b3:cell-cycle-apoptosis:checkpoints}{जाँच-बिंदु} ने Cdc25 को रोककर S-वक्र दाईं ओर खिसका दिया है, इसलिए
छलाँग नहीं लगती। मरम्मत पर Cdc25 मुक्त होता है, देहली जमा हुए \omterm{def:b3:cell-cycle-apoptosis:cdk}{चक्रिन} से
नीचे गिर जाती है, और कोशिका तुरंत समसूत्री विभाजन में घुस जाती है ---
इसीलिए \omterm{def:b3:cell-cycle-apoptosis:checkpoints}{जाँच-बिंदु} से छूटी कोशिकाएँ एक साथ समसूत्री विभाजन में जाती हैं।
\textbf{12.} न \omterm{def:b3:cell-cycle-apoptosis:cdk}{चक्रिन} B अपघटित होता है न सेक्युरिन: कोशिका कोहेसिन अक्षत
और Cdk1 सक्रिय रखे अनिश्चित काल तक मध्यावस्था में पड़ी रहती है; वह वहीं
मर जाती है या अंततः किसी अविभाजित जीनोम के साथ फिसलकर निकल जाती है।
\textbf{13.} प्रति दिन $10^{11}$ विभाजन।
\textbf{14.} \omterm{def:b3:cell-cycle-apoptosis:checkpoints}{जाँच-बिंदु} के साथ प्रति दिन $10^{11}/5\times 10^{4} = 2\times 10^{6}$
असुगुणित पुत्रियाँ; उसके बिना $2\times 10^{9}$।
\textbf{15.} प्रति दिन $2\times 10^{4}$; $80$ वर्षों में $6\times 10^{8}$।
\textbf{16.} प्रति दिन $10^{9}/50 = 2\times 10^{7}$ गलत बँटवारे: हर
दिन अर्बुद दो करोड़ नए गुणसूत्र-प्ररूप आज़माता है, और चयन उन्हें रख
लेता है जो तेज़ी से बढ़ें या किसी औषध का प्रतिरोध करें।
\textbf{17.} हर विभाजन में गलत बँटवारा होने पर कोई भी भ्रूणीय
कोशिका-वंशक्रम सुगुणित जीनोम नहीं रख पाता और भ्रूण मर जाता है। आंशिक
ह्रास उन कोशिकाओं में कभी-कभार असुगुणिता देता है जो उसे झेल जाएँ
(विशेषकर p53 के बिना), यानी कोई गुणसूत्री अस्थिरता जो जीव को मारे बिना
अर्बुद को विचरण देती है।
\textbf{18.} कोई अर्धसूत्री त्रुटि अतिरिक्त गुणसूत्र को युग्मक में डाल
देती है, इसलिए युग्मनज और उससे बनी हर कोशिका त्रिसूत्री होती है; कोई
समसूत्री त्रुटि किसी एक कायिक कोशिका में होती है और केवल उसके क्लोन को
वंशागत होती है।
\textbf{19.} प्रति दिन $10^{11}\times \qty{1}{ng} = \qty{100}{g}$,
प्रति वर्ष \qty{36}{kg} --- हर वर्ष लगभग शरीर के आधे द्रव्यमान जितना।
\textbf{20.} प्रति दिन \qty{20}{g} प्रोटीन, यानी आहारी सेवन का लगभग
एक-तिहाई और शरीर के कुल प्रोटीन-आवागमन का कुछ प्रतिशत।
\textbf{21.} $10^{11}$ मुर्दे, प्रति महाभक्षक प्रति दिन $24$ की दर से:
पूरे समय काम करते $4\times
10^{9}$ महाभक्षक, यानी $10^{11}$ महाभक्षकों के समय का \qty{4}{\%}।
\textbf{22.} कोई लयित मुर्दा प्रोटिएज़, DNA, ATP और दूसरे
चेतावनी-संकेत छोड़ता है जो शोथ करते हैं और केंद्रकीय प्रतिजनों के
विरुद्ध स्व-प्रतिरक्षा भड़का सकते हैं; अक्षत मुर्दा अपने बाहरी परत पर
\omterm{def:b3:cell-cycle-apoptosis:apoptosis}{फ़ॉस्फ़ेटिडिलसेरीन} से अपना विज्ञापन करता है (और अपने ``मुझे मत खाओ''
संकेत खोकर), और मुक्त न्यूक्लियोटाइडों से भक्षकाणुओं को बुलाता है।
\textbf{23.} \omterm{def:b3:cell-cycle-apoptosis:pathways}{आंतरिक पथ} सूत्रकणिका पर रुका है। बाह्य मार्ग वहाँ अब भी
काम करता है जहाँ कैस्पेज़-8 कैस्पेज़-3 को सीधे सक्रिय करती है, यद्यपि
Bid के ज़रिए उसका प्रवर्धन खो जाता है। कोशिकाएँ तेज़ी से विभाजित होकर
नहीं, न मरकर जमा होती हैं --- धीमी वृद्धि --- जब तक कोई बाद का घाव
(प्रायः \emph{MYC}) प्रचुरोद्भवन न जोड़ दे।
\textbf{24.} \omterm{def:b3:cell-cycle-apoptosis:apoptosis}{एपॉप्टोसिस} घंटों लेता है, उसे ATP चाहिए और वह ऐसे चरणों से
गुज़रता है जिन्हें रोका जा सकता है --- \omterm{def:b3:cell-cycle-apoptosis:apoptosis}{कैस्पेज़} निरोधक, सूत्रकणिकीय
देहली --- इसलिए समय पर उपचार मिलने पर अर्ध-क्षेत्र के न्यूरॉन बचाए जा
सकते हैं; क्रोड के न्यूरॉन मिनटों में ऊर्जा की विफलता से मरते हैं, और
बीच में रोकने को कोई कार्यक्रम ही नहीं होता।
\textbf{25.} नंगे दोलक का आवर्तकाल लगभग \qty{16}{h}; किसी सामान्य शरीर
में प्रति दिन कोई $2\times
10^{4}$ असुगुणित विभाजनशील कोशिकाएँ; \omterm{def:b3:cell-cycle-apoptosis:apoptosis}{एपॉप्टोसिस} से हर दिन \qty{100}{g}
कोशिकाएँ पुनर्चक्रित।
\end{solution}
